\section{Đếm phương tiện}

\begin{frame}
\frametitle{Bài toán đếm phương tiện}
\framesubtitle{Từ quỹ đạo theo dõi đến thống kê lưu lượng}

\begin{itemize}[<+->]
  \item Đầu vào là các vết theo dõi đã có ID ổn định từ giai đoạn tracking
  \item Mỗi phương tiện được biểu diễn bởi hộp giới hạn, ID và lớp phương tiện
  \item Hệ thống kiểm tra vị trí đại diện của phương tiện với vùng đếm đã cấu
        hình
  \item Mỗi ID chỉ được cộng một lần trong từng vùng đếm để tránh đếm trùng
  \item Kết quả gồm số lượng hiện thời, tổng tích lũy và thống kê theo lớp
\end{itemize}

\vspace{0.2cm}

{\scriptsize Cách tiếp cận dựa trên quỹ đạo phù hợp với hướng
detection-based counting \cite{dai2019,shokri2025}.}

\end{frame}

\begin{frame}
\frametitle{Vùng đếm}
\framesubtitle{ROI dạng đa giác}

\begin{itemize}[<+->]
  \item Người dùng tạo một hoặc nhiều vùng quan tâm (ROI) trực tiếp trên khung
        hình video
  \item ROI được biểu diễn bằng đa giác để phù hợp với phối cảnh camera giao
        thông
  \item Điểm đại diện của phương tiện được chọn tại tâm đáy hộp giới hạn
  \item Nếu điểm đại diện nằm trong ROI, phương tiện được xem là đang thuộc
        vùng đó
  \item Cho phép sửa vị trí đỉnh ROI và xóa ROI để hiệu chỉnh theo từng video
\end{itemize}

\end{frame}

\begin{frame}
\frametitle{Thuật toán đếm}
\framesubtitle{Counting by tracked IDs}

\begin{enumerate}
  \item Nhận danh sách track hợp lệ ở khung hình hiện tại
  \item Với mỗi track, lấy ID, lớp phương tiện và hộp giới hạn
  \item Tính điểm đại diện tại tâm đáy hộp giới hạn
  \item Kiểm tra điểm đại diện có nằm trong ROI nào hay không
  \item Nếu nằm trong ROI, tăng bộ đếm live và bộ đếm theo lớp
  \item Nếu ID chưa từng được đếm trong ROI đó, tăng tổng tích lũy
\end{enumerate}

\vspace{0.25cm}

\begin{block}{Nguyên tắc chống đếm trùng}
  Mỗi ROI lưu tập ID đã được ghi nhận. Một phương tiện chỉ làm tăng tổng tích
  lũy lần đầu tiên khi ID của nó xuất hiện trong ROI.
\end{block}

\end{frame}

\begin{frame}
\frametitle{Công cụ minh họa}
\framesubtitle{Interactive tracking counter}

\begin{itemize}[<+->]
  \item Xây dựng script chạy độc lập từ dòng lệnh:
        \texttt{src/run\_tracking\_counter.py}
  \item Hỗ trợ cấu hình riêng qua \texttt{.apprc} và vẫn dùng cấu hình dữ liệu
        chung của hệ thống
  \item Cho phép chọn tracker, thiết bị chạy, ngưỡng phát hiện và đường dẫn
        video từ CLI
  \item Hỗ trợ tạo, chọn, sửa và xóa ROI trong cửa sổ hiển thị
  \item Có thể lưu video kết quả bằng \texttt{--save-output} hoặc
        \texttt{--output}
\end{itemize}

\end{frame}

\begin{frame}
\frametitle{Bảng theo dõi và đếm}
\framesubtitle{Thông tin hiển thị khi chạy}

\begin{itemize}[<+->]
  \item Hiển thị FPS và tổng số ID duy nhất đã được ghi nhận
  \item Với mỗi ROI, hiển thị số lượng phương tiện đang nằm trong vùng và tổng
        số đã đếm
  \item Thống kê nhanh số lượng theo lớp phương tiện như car, bus, van hoặc
        others
  \item Hộp giới hạn hiển thị ID, lớp phương tiện và hướng di chuyển lên/xuống
  \item Có thể bật/tắt bảng điều khiển và bảng counting bằng phím tắt trong
        quá trình chạy
\end{itemize}

\end{frame}

\begin{frame}
\frametitle{Nhận xét}
\framesubtitle{Ưu điểm và hạn chế}

\begin{itemize}[<+->]
  \item Cách đếm đơn giản, trực quan và dễ áp dụng cho các camera cố định
  \item Việc dùng ID theo dõi giúp hạn chế đếm trùng khi phương tiện xuất hiện
        nhiều khung hình liên tiếp
  \item ROI dạng đa giác linh hoạt hơn đường đếm cố định trong các góc nhìn
        phối cảnh
  \item Kết quả đếm phụ thuộc mạnh vào chất lượng phát hiện và độ ổn định ID
        của tracker
  \item Cần đánh giá định lượng sâu hơn trên các kịch bản đếm thực tế và các
        vùng đếm được gán nhãn chuẩn
\end{itemize}

\end{frame}
